% ============================================================
%  第一章  集合与常用逻辑用语
%  1.1  集合的概念
% ============================================================

\section{集合的概念}

\subsection{知识点梳理}

\begin{definition}[集合与元素]
一般地，我们把研究对象统称为\emph{元素}，把一些元素组成的总体叫做\emph{集合}，简称为集．通常用大写拉丁字母 $A,B,C,\dots$ 表示集合，用小写拉丁字母 $a,b,c,\dots$ 表示元素．

如果 $a$ 是集合 $A$ 的元素，就说 $a$ \emph{属于}集合 $A$，记作 $a\in A$；如果 $a$ 不是集合 $A$ 的元素，就说 $a$ \emph{不属于}集合 $A$，记作 $a\notin A$．
\end{definition}

\begin{definition}[集合中元素的特性]
\begin{enumerate}
  \item \emph{确定性}：给定一个集合，任何一个对象是不是这个集合的元素是确定的．
  \item \emph{互异性}：集合中的元素互不相同，同一个元素在集合中只能出现一次．
  \item \emph{无序性}：集合中的元素没有先后顺序．
\end{enumerate}
\end{definition}

\begin{note}[常用数集及其记法]
非负整数集（自然数集）记作 $\N$；正整数集记作 $\N^{*}$ 或 $\N_{+}$；整数集记作 $\Z$；有理数集记作 $\Q$；实数集记作 $\R$．另外，不含任何元素的集合叫做\emph{空集}，记作 $\varnothing$．
\end{note}

\subsection{集合的表示方法}

\begin{definition}[列举法与描述法]
\begin{enumerate}
  \item \emph{列举法}：把集合的所有元素一一列举出来，并用花括号“$\{\ \}$”括起来表示集合的方法，如 $\set{1,2,3}$．
  \item \emph{描述法}：用集合中元素的共同特征表示集合的方法，形如 $\setof{x\in A}{p(x)}$，其中 $p(x)$ 为元素 $x$ 所具有的共同特征，如 $\setof{x\in\R}{x>4}$．
\end{enumerate}
\end{definition}

\subsection{典型例题}

\begin{example}{元素与集合的关系}
给出下列关系：① $\dfrac{1}{2}\in\R$；② $\sqrt{2}\notin\Q$；③ $\abs{-3}\in\N^{*}$；④ $0\in\varnothing$．其中正确的个数是（\quad）

\begin{choices}
  \item $1$ \item $2$ \item $3$ \item $4$
\end{choices}

\begin{solution}
① 因为 $\dfrac{1}{2}$ 是实数，所以 $\dfrac{1}{2}\in\R$ 正确；

② 因为 $\sqrt{2}$ 是无理数，不是有理数，所以 $\sqrt{2}\notin\Q$ 正确；

③ 因为 $\abs{-3}=3$ 是正整数，所以 $\abs{-3}\in\N^{*}$ 正确；

④ 空集中不含任何元素，所以 $0\in\varnothing$ 错误．

故正确的个数为 $3$，选 C．
\end{solution}
\end{example}

\begin{example}{元素的互异性}
已知集合 $A=\set{a+2,\,(a+1)^{2},\,a^{2}+3a+3}$，若 $1\in A$，求实数 $a$ 的值．

\begin{analysis}
$1\in A$ 说明 $A$ 的三个元素中至少有一个等于 $1$，故需分三种情况讨论；求出 $a$ 后还必须检验集合中的三个元素是否互不相同．
\end{analysis}

\begin{solution}
由 $1\in A$，分三种情况：

（1）若 $a+2=1$，则 $a=-1$，此时 $A=\set{1,0,1}$，元素 $1$ 重复，不符合互异性，舍去；

（2）若 $(a+1)^{2}=1$，则 $a=0$ 或 $a=-2$．

当 $a=0$ 时，$A=\set{2,1,3}$，符合题意；

当 $a=-2$ 时，$A=\set{0,1,1}$，元素 $1$ 重复，舍去；

（3）若 $a^{2}+3a+3=1$，即 $a^{2}+3a+2=0$，解得 $a=-1$ 或 $a=-2$，由（1）（2）知均不符合题意．

综上，$a=0$．
\end{solution}
\end{example}

\begin{example}{集合的表示方法}
用适当的方法表示下列集合：

（1）方程 $x^{2}-2x+1=0$ 的所有实数根组成的集合；

（2）大于 $1$ 且小于 $6$ 的所有整数组成的集合；

（3）不等式 $2x-3>5$ 的解集．

\begin{solution}
（1）方程 $x^{2}-2x+1=0$ 即 $(x-1)^{2}=0$，解得 $x=1$，故所求集合为 $\set{1}$．

（2）大于 $1$ 且小于 $6$ 的整数为 $2,3,4,5$，故所求集合为 $\set{2,3,4,5}$．

（3）由 $2x-3>5$ 得 $x>4$，故所求集合为 $\setof{x\in\R}{x>4}$．
\end{solution}
\end{example}

\begin{method}[判断能否构成集合]
判断一组对象能否构成集合，关键看是否满足\key{确定性}：对象的标准是否明确．“著名”“接近”“比较小”这类描述没有明确标准，不能构成集合；而“方程的解”“某班身高超过 170 cm 的同学”标准明确，可以构成集合．
\end{method}

\begin{note}[三个易错点]
\begin{enumerate}
  \item 集合 $\set{1,1,2}$ 与 $\set{1,2}$ 表示同一个集合，列举时不可重复书写元素．
  \item $\setof{x}{y=x^{2}}$ 与 $\setof{y}{y=x^{2}}$ 一般是不同的集合：前者是函数自变量的取值集合（定义域），后者是函数值的集合（值域）．
  \item $0$、$\set{0}$、$\varnothing$ 三者不同：$0$ 是一个数，$\set{0}$ 是含有一个元素的集合，$\varnothing$ 是不含任何元素的集合，且 $\varnothing\ne\set{0}$．
\end{enumerate}
\end{note}

\subsection{课堂练习}

\begin{exercise}{基础巩固}
下列各组对象能构成集合的是（\quad）

\begin{choices}
  \item 我国所有的大河
  \item 接近 $0$ 的全体实数
  \item 方程 $x^{2}-3x+2=0$ 的实数根
  \item 我国著名的数学家
\end{choices}
\end{exercise}

\begin{exercise}{基础巩固}
用列举法表示集合 $B=\setof{x\in\N}{1<x<5}$．
\end{exercise}

\begin{exercise}{能力提升}
已知集合 $M=\set{1,\,a,\,a^{2}-a}$，若 $3\in M$，求实数 $a$ 的值．
\end{exercise}

\begin{exercise}{能力提升}
已知集合 $A=\setof{x\in\R}{ax^{2}+2x+1=0}$，若 $A$ 中只有一个元素，求实数 $a$ 的值．
\end{exercise}
